\chapter{Fazit}
\begin{itemize}
    \item Noch Bumper, Abstandssensoren und Problem mit Codebasis ansprechen(Alte Leichen in den GVLs)
    \item Probleme mit anderen Gewerken ansprechen (Kommunikation, verlässlichkeit, Hebel)
    \item Arbeit über Semesterferien früher kommunizieren. War am Anfang nicht klar.
    \item Grundsätzlich aber viel mitgenommen, Freiheiten sehr genossen.
    \item Kein Plichtenheft
    \item Kein Projektleiter
    
    \text{Regler}
    Die Reglerstruktur und die Parameter sind gut geeignet um die Anforderungen an das Projekt zu erfüllen und um Gewerk 2 in die Lage zu versetzen eine sichere Bahnfahrt zu realisieren. Die Reglerstruktur war von Beginn des Projektes unser erster Fokus gewesen und wurde bereits im Januar, noch ohne Beobachter realisiert, um den anderen Gewerken das Testen zu ermöglichen. Der Beobachter hat im weitern Verlauf die Performance enorm verbessert, ebenso wie das Anpassen der Regelparameter. Im Kontext des Projekts sind wir sehr zufrieden mit dem Regler. Obwohl es für die Anforderungen des Projekts nicht notwendig gewesen wäre, hätten wir uns gerne etwas mit quantitativen Gütekriterien für den Regler beschäftigt, um die Performance weiter zu verbessern. Dafür ist ein lauffähiger, gleichbleibender Stand der Roboterfabrik nötig. Leider konnte Gewerk 2 diesen Stand bis kurz vor Projektende nicht präsentieren, weshalb die quantitative Regleroptimierung nicht stattgefunden hat.
\end{itemize}